\chapter[Strangers Inside Us: Novels]{How Novels Let Strangers Live Inside Us}
\section{An Old Book Full of Strangers' Notes}
Pick something up, not new this time: an old book.

Suppose you bought it for five yuan at a weekend stall: a dog-eared novel with a fraying cover and two strips of tape around its spine. Opening it reveals several previous tenants, their pencil notes crowding the margins. A neat hand records on the flyleaf, "Bought at the county Xinhua Bookshop, 1987." Halfway through, another reader heavily underlines dialogue in ballpoint and writes "Nonsense!" Several pages later, pale pencil challenges that judgement: "What nonsense? I think it's true." Near the back, one page is deeply wrinkled and stiff, as if wetted then dried. Its margin has no words, only three emphatic exclamation marks.

Holding it, you realise these strangers met on one page across decades. The character who made them underline, argue, and cry never existed. The conversation that wrinkled the page was invented beneath a writer's desk lamp.

This is our question: why can black marks on white paper, openly announcing their fiction, move generations of strangers? How do novels bring people who have never met, including nonexistent people, into our hearts?

\section{A Jianyang Bookshop in the Wanli Reign}
Come to another scene: late-sixteenth- or early-seventeenth-century Ming China, during the Wanli reign, among the publishers of Masha and Chonghua in Jianyang, Fujian.

Before reaching the street, you smell ink, freshly split pearwood, and paste. Bookshops line both sides, their fronts opening onto workshops. Inside, carvers bend over low tables. Thin handwritten sheets lie reversed on wooden blocks, and knives make small seed-cracking sounds. A skilled worker carves perhaps a hundred characters daily. A novel of hundreds of thousands of characters requires hundreds or thousands of blocks, filling half a room.

In the courtyard, printers lay paper on inked blocks, press with palm-fibre brushes, and lift finished pages. Drying sheets hang on lines, making the courtyard's words flutter in the wind. At the front counter, a bookseller bargains with distant buyers over wholesale prices for a newly cut historical romance and shipping times to Nanjing or Suzhou.

The book being rushed surprises you. Not classics or examination essays, but a long illustrated story of wars and outlaw heroes. Bound volumes on the counter offer gods and demons, criminal cases, talented scholars and beautiful women. Business is excellent, the owner says. Buyers include merchants, clerks, frustrated examination candidates, officers at leisure, and literate women purchasing through others. They read neither for examinations nor moral cultivation, but enjoyment.

Remember this workshop: the ingredients of our later arguments all inhabit its courtyard. Printing cheapens words, cities and markets support leisure reading, and ordinary readers can reach books. Long fiction grew from this bustle, not by dropping from a scholar's study.

\section[Fictional Events, Real Tears]{Knowing It Is Fiction, Why Are the Tears Real?}
Set out the conflict before answering.

There are two forceful attitudes towards crying over novels.

One is dismissive: do we have so many tears for real people that we can spare some for invented ones? A page's character has no breath or birthday and cannot know you cried until midnight. Call it sentimentality kindly, self-deception unkindly: printed tricks remotely control your feelings. Why be proud? Better care about someone genuinely hungry downstairs.

The opposite view finds the tears extraordinary precisely because their object does not exist. Neighbours share daily encounters; relatives share blood and interests. What do you gain by grieving for someone wholly fictional and unrelated? Nothing. That reveals a human capacity to cross kinship, geography, even existence itself, and feel another life's circumstances. Novels train this capacity.

Both have a point, and both lead beyond whether we ought to cry:

\textbf{How does fiction produce real feeling?}

Fiction advertises itself across cultures and centuries. Storytellers begin "The story goes"; Western covers announce "A novel". Authors never meant to claim factual truth, and readers were never fooled. Yet acknowledged invention unleashes emotion. This resembles conspiracy rather than deception: the writer builds a stage, readers enter voluntarily, invite an unreal person inside, then worry and lose sleep for them.

How was this stage built, and who invented the machinery? It has held our best and worst impulses. Before continuing, locate yourself: do tears for fictional people count?

\section[Officials, Booksellers, and Prosecutors]{Minor Officials, Booksellers, and Prosecutors}
At least four historical voices speak about these leisure books. Hear them all.

The oldest is orthodoxy. Xiaoshuo, today's word for fiction, first appears in the Zhuangzi's "External Things": ornamenting petty talk to seek distinction remains far from great understanding. Here xiaoshuo means trivial discourse, not modern novels, but its dismissive tone marks a lowly birth. In the Eastern Han, Ban Gu's bibliographical treatise in the Book of Han ranks this school last:

\begin{quote}
The school of minor tales probably arose among petty officials, making its material from street talk, alley gossip, and things heard along the road.
\end{quote}
Baiguan were minor officials collecting popular talk. In effect, fiction descended from imperial gossip-gatherers dealing in hearsay. Ban Gu adds that gentlemen disdain it, yet ordinary people enjoy it, so it cannot disappear. This is an early love-hate assessment: contempt coupled with recognition of survival.

The second voice belongs to markets and readers voting with coins and tears. Jianyang demonstrates it: publishers compete to print what orthodox scholars despise, and readers buy. The eccentric late-Ming, early-Qing critic Jin Shengtan provocatively grouped Zhuangzi, "Encountering Sorrow", Records of the Grand Historian, Du Fu's poems, Water Margin, and Romance of the Western Chamber as Six Works of Genius. He elevated an outlaw novel alongside the great history, annotating it chapter by chapter and praising its craft. What critics called gossip, he deliberately called genius.

The third voice is government, internationally recognisable. Ming and Qing rulers and local officials repeatedly banned and destroyed licentious fiction for encouraging lust, theft, and moral decay. France prosecuted Flaubert in 1857 over Madame Bovary, the story of a provincial doctor's wife, marital discontent, adultery, and destruction. He was acquitted, but the prosecutor's concern deserves remembering: readers, especially women, might regard indulgence as understandable or even sympathetic.

Do not laugh at censors too quickly. Give their case fully, a necessary exercise when their position is usually treated as ridiculous.

Their argument has three layers. First, imitation is real. Goethe's epistolary Sorrows of Young Werther, published in 1774, follows a young man loving an engaged woman and finally shooting himself. Europe imitated his clothes, and reports described imitative suicides. Sociologists later named imitation following suicide coverage the Werther effect. Fears of readers copying stories were not wholly invented. Second, social order is real. Where parents arrange marriages and fixed status organises society, sympathy for adulterers or outlaws loosens structural screws. Third, protection feels real to censors, who see themselves guarding those deemed vulnerable: youth, women, and people without classical schooling.

The costs must also be clear. Who decides vulnerability? Why assume women and ordinary people deserve only filtered worlds? Giving officials the gatekeeping key declares some people permanently unfit to grow up. Bans also often backfire: notoriety and prices rise first. Ming and Qing prohibited-book lists now resemble recommendations of their periods' most enjoyable reading.

The fourth voice, often drowned out, is ordinary readers. They neither publish nor petition, but undeniably exist. Contemporary Edo, now Tokyo, had kashihonya rental bookshops, hundreds in the city at their height. Workers carried boxes of books through streets; samurai wives, shopkeepers, and artisans cheaply borrowed new popular fiction for days. Edo's literacy was remarkably high for its time. These readers left no famous pronouncements or monuments, merely admitted strangers' stories into their homes night after night.

Contempt, praise, fear, dependence: fiction has lived at their contested centre from its beginning.

\section[Thinking Bridge: Writer, Voice, View]{A Thinking Bridge: Who Wrote, Who Speaks, Who Sees?}
The first reading tool separates three questions: \textbf{who wrote this, who tells it, and whose eyes are we looking through?} They often have different answers. Confusing them creates accusations that authors defend murderers and leaves readers vulnerable to skilful narration.

\begin{itemize}
\item \textbf{Who wrote it}: The flesh-and-blood author, who eats and receives royalties.
\item \textbf{Who tells it}: The narrator, a voice the author invents. A first-person "I" is usually a character, not the writer.
\item \textbf{Who sees}: The viewpoint followed at this moment. One room becomes two different rooms through a child's eyes and a butler's.
\end{itemize}
These questions reveal something elegant: \textbf{centuries of novelistic development continually rearrange these three positions}. Consider four major arrangements.

\textbf{Chaptered storytelling: the storyteller presides.} Romance of the Three Kingdoms and Water Margin have invisible, omnipresent narrators opening "The story goes" and closing "To learn what happened next, listen to the next instalment". They know commanders' thoughts and simultaneous events far away, interrupt with opinions, and create suspense. Real teahouse storytelling produced this form: tomorrow's tea money required a hook today. The printed closing formula fossilises the spoken one.

\textbf{Epistolary fiction: the author withdraws and readers open private letters.} Eighteenth-century Europe made novels from correspondence. Richardson's Pamela, 1740, consists of a maid's letters to her parents about her master's advances and her resistance. Goethe's Werther, 1774, is mostly letters to a friend. The author vanishes, characters narrate, and readers become intruders opening correspondence never intended for them, intensifying apparent authenticity. Its limitation creates its fascination: we know only what writers choose and manage to express. Through a self-uncomprehending person's words, we must infer the person.

\textbf{Realism: an invisible all-seeing camera.} Nineteenth-century French and Russian novels move narrators backstage, unnamed and unseen yet knowing every drawing-room exchange and every unspoken thought. Balzac's more than ninety works in The Human Comedy position him as French society's secretary, registering an age. Flaubert seeks an author like God, everywhere present but never speaking. This invisible reliable narration remains our default image of a novel.

\textbf{Modernism: remove the camera and enter the mind.} Early-twentieth-century Joyce, Woolf, and Faulkner found the camera too external. Thoughts are not fluent sentences but currents of smells, memories, fragments, and heartbeat. Could novels write that flow directly? Stream of consciousness borrows William James's description of continuous thought in his 1890 Principles of Psychology. An ordinary Dubliner's day in Ulysses becomes hundreds of pages of mental current. In Faulkner's Sound and the Fury, brothers narrate one family differently, their worlds combining towards truth.

These are divisions of labour, not ranks: storytellers entertain, letters bring intimacy, cameras breadth, consciousness depth. All discover that interior life deserves writing, hundreds of pages for otherwise unspeakable thoughts. Almost no genre had done this before the novel matured.

Now remove a widespread misconception: \textbf{realism is not an exact photocopy of reality}.

Stendhal famously likens the novel in The Red and the Black to a mirror carried along a road, reflecting sky or mud. But mirrors do not move themselves. Someone chooses the road, angle, and time spent before each puddle. Faithful portrayal contains thousands of decisions: whom to include, where to begin, which speech emerges or is swallowed. Dream of the Red Chamber's crab feast reveals hierarchy, talent, and private troubles because Cao Xueqin selected and constructed it from countless possible dinners. Realism is a crafted effect, not a photocopy. Remembering this prevents direct treatment of novels as historical evidence and prompts a better question when something feels utterly true: who made it feel so, and how?

\section[Who Understands These Absurd Words?]{Pages of Absurd Words: Who Understands Their Flavour?}
Read a little primary evidence. China's most admired novel opens with its author reflecting on fiction.

After explaining the book's origins in Dream of the Red Chamber's first chapter, Cao Xueqin adds:

\begin{quote}
Pages full of absurd words, a handful of bitter tears! Everyone calls the author foolish; who understands their flavour?
\end{quote}
At the entrance to a dream realm in the same chapter hangs a couplet:

\begin{quote}
When false is taken for true, true too becomes false; where nothing is made something, something returns to nothing.
\end{quote}
Another opening, the first sentence of Romance of the Three Kingdoms, declares:

\begin{quote}
The story goes that the world's great affairs, long divided, must unite; long united, must divide.
\end{quote}
Together these passages nearly map Chinese fiction's whole attitude to invention.

First, huashuo, "the story goes", openly reveals the storyteller's origin. The narrator assumes listeners like a teahouse audience. More fundamentally, this is not presented as news. A contract is openly signed: I tell a story; you receive it as one.

The couplet then says that cultivating falsehood as truth can expose apparent reality as false. Read it as method: serious fictional invention reveals how honours, respectability, and family glory may be more theatrical than novels. Fiction is not reality's opposite but its developing fluid.

The poem voices the author's side of our question. Its contradiction matters: he admits absurd words without defence, yet answers with bitter tears, real enough to gather in handfuls. "Who understands their flavour?" is sadder still. He knowingly invents and fears nobody will understand the truth inside invention. This and the old book's exclamation marks are opposite ends of one exchange: the writer fears being misunderstood, the reader urgently says they understand.

We did not invent this doorway between fiction and truth. People have passed through it for centuries.

\section{Why Did Novels Grow Up Just Then?}
Across the map, popular long narratives rose from low-status curiosities to major cultural commodities: Ming--Qing China, Edo Japan, eighteenth- and nineteenth-century Europe. Dates differ, patterns resemble one another. Why? Evidence does not interpret itself. Compare three explanations, each with strengths and limits.

\textbf{Explanation 1: Technology made books cheaper.}

Manuscript labour cost more than ordinary households could afford; printing reduced reproduction costs dramatically. Chinese woodblocks and Gutenberg's movable type turned luxuries into commodities. Jianyang's production line forms part of that chain. Its strength is identifying material possibility: mass reading needs affordable books. Its weakness is timing. Chinese printing flourished in the Song, yet chaptered fiction peaked in the later Ming. Europe had type in the fifteenth century but waited three centuries for the novel's boom. Technology permits; it does not cause occurrence. Cheap cooking pots alone do not open restaurants.

\textbf{Explanation 2: Cities, markets, and literate readers came together.}

This supplies demand. Ming--Qing Jiangnan, Edo Japan, and eighteenth-century England shared expanding commercial cities and literate people with spare money: merchants, assistants, artisans, household women. Reading offered leisure and others' lives, not official memorial-writing. English circulating libraries rented mostly novels; Edo's book boxes and Chinese booksellers' boats served comparable customers. This solidly explains who bought and how much, but not form. Opera, gambling, and music also entertained. Why did long novels catch these readers?

\textbf{Explanation 3: The individual emerged and needed a room large enough.}

This goes inward. Epics concern heroes and gods, drama scenes and songs, histories rulers and generals. Ordinary city people, merchants' sons and doctors' wives, began feeling their whole lives deserved telling. Modern society loosened fixed kinship and status, making individuality an experience: my love, future, and grievances matter. The novel uniquely accommodates it, following ordinary lives across decades, entering thoughts, exploring a dinner's unspoken concerns. This explains protagonists moving from heroes to ordinary people. But cause and effect may reverse. Did modern individuals precede novels, or did novels help create them? Probably both: readers learnt to inspect themselves like characters, multiplying their inner dramas.

Long stories need not follow only printing plus cities. The Arabian Nights accumulated over centuries from Persian foundations, Baghdad tales, and Cairo anecdotes, through intertwined oral and manuscript circulation. Scheherazade's nightly stories to preserve her life join hundreds into one collection. Its growth did not depend on printing. Human appetite for long stories predates presses; presses merely let it eat its fill.

Materials, markets, minds: no explanation replaces another. Taking one alone as truth flattens history.

\section[Further Challenge: Polyphony]{A Deeper Challenge: Polyphony Lets Characters Live Beyond Their Author}
Modernism took novels deeper inside one mind. Our strongest theoretical question points the opposite way: how many equal minds can one novel contain?

Soviet theorist Bakhtin, 1895--1975, posed it while studying Dostoevsky, author of Crime and Punishment and The Brothers Karamazov. His Problems of Dostoevsky's Poetics introduced \textbf{the polyphonic novel}.

Roughly, most earlier novels, including great ones, were monologic: many characters ultimately spoke lines assigned by one authorial voice. The author had already judged their arguments, and characters enacted the verdict like puppets. Dostoevsky revolutionised this by giving characters complete, independent voices. A murderer debates the author as an equal; an atheist silences a believer without authorial adjudication. Conflicting truths coexist and collide without any fully dominating, like multiple independent musical lines woven together rather than one melody with accompaniment. That is polyphony, many-voicedness.

Apply this measure to Dream of the Red Chamber. Jia Baoyu has complete reasons to reject worldly advancement; Xue Baochai has hers for urging it; Wang Xifeng's calculations make sense within her position. Cao Xueqin depicts arguments and mistakes but refuses an official final stance. Centuries of fierce Daiyu-versus-Baochai debate persist because the author withholds judgement. \textbf{Polyphony offers no standard answer but several internally persuasive worlds through which readers must travel}.

Why is that remarkable? It trains moral capacity. Monologic fiction rehearses choosing the right side; polyphony asks something harder: inhabit a mind you disagree with, see how coherent its reasons feel from within, then emerge still carrying some of its logic. It extends our practice of giving opponents their fullest case across hundreds of pages.

Mark the theory's boundaries before it becomes prejudice. First, not all good novels are polyphonic. Fables, satire, and fairy tales often gain strength from simplicity; Journey to the West need not let demons debate its author. Polyphony is a method, not a medal. Second, silence is not authorial absence. Choosing speakers, eloquence, and silences remains power. Polyphony turns the author from judge into parliamentary chair: no direct verdict, but control of the agenda. Recognising that helps resist apparent neutrality in fiction, news, and short videos alike.

\section[The Bookstall in the Algorithmic Age]{The Second-Hand Bookstall in the Algorithmic Age}
This centuries-old issue lives in your phone.

Online novels post thousands of words daily while readers demand updates: "listen next time" with fresh batteries. Wish-fulfilment fiction times emotional revenge as storytellers timed striking their wooden blocks. Fan fiction turns marginal notes into new main texts, readers borrowing characters to continue them. Review sections and scrolling comments enlarge pencilled margins into public squares. Strangers still meet, argue, and correct one another in stories' spaces, now seconds apart rather than decades.

Recommendation algorithms are the Jianyang bookseller's strongest descendants, remembering your tastes better than you and filling your arms around the clock with suggested favourites.

Face the central controversy: \textbf{does reading novels actually make people kinder?}

Optimists have many witnesses. Uncle Tom's Cabin appeared in 1852, sold hundreds of thousands of copies within a year, introduced Northern readers to an enslaved person's life, and strengthened abolitionist feeling. Lincoln reportedly called Harriet Beecher Stowe the little woman who caused the great war, but reliable documentation is absent: treat it as legend. Its survival nevertheless records belief in fiction's power. Psychology offered support when a 2013 study reported short-term improvements in understanding other minds after literary fiction. Honesty requires adding that later replication attempts produced mixed outcomes and no settled consensus. Plausibly helpful is not guaranteed helpful.

A counterexample warns against concluding that more reading means greater kindness. History provides no guarantee. Twentieth-century killing machines included cultivated people who loved Goethe and wept for Beethoven, cultured executioners in concentration camps. Crying for characters is safe practice: they neither reject sympathy nor demand anything afterwards. Real compassion costs. Mistaking goodness on the page for goodness in life is a reader's most respectable trap.

Fiction can also reinforce prejudice. Recurring stereotypes make particular regions calculating, professions cruel, or women exist only to be rescued. Story carries these images into hearts alongside sympathy, harder to remove than textbook prejudice because they feel personally experienced rather than taught. Who receives interior life and who becomes scenery always matters. Giving someone a whole chapter of thought implicitly licenses whose feelings deserve attention.

The question today is not whether to read but with which tools. Separate writer, narrator, and viewpoint; count genuinely independent voices; ask how reality effects are made; distrust effortless comfort, which may fit your existing prejudices exactly.

\section[Your Turn: Tears for Fictional People]{Your Turn: Tears for Someone Who Never Existed}
Return to the old bookstall volume.

Its tear-wrinkled page, three exclamation marks, and reader rebuked by a later stranger are real; the character causing them is not. Now answer with reasons, though no standard answer exists:

\textbf{Do tears for fictional characters make you better, or merely make you feel good about yourself?}

Weigh both sides.

One says the tear genuinely reaches another life without request or reward. Applied outside fiction, that capacity becomes compassion. Perhaps every feeling that strangers' affairs concern us needs safe rehearsal first, and novels provide the rehearsal room.

The other calls it cheap consumption. A few leisure hours purchase a flattering self-image. Characters will not remain hungry after the book closes or need money, advocacy, and uncomfortable action. You can rehearse expertly forever without going onstage.

Go deeper: could both be right? Is there a way to turn tears into steps rather than a sofa? Has a particular book actually changed your attitude towards real people?

After decades of passing from hand to hand, the book reaches you. One line remains empty in the margin. What will you write?
